\summarypart{EIGHT}{The Age of Ideologies}{1914--1989}
\begin{partoverview}
The twentieth century brought the colossal struggles of political and military giants: Communism,
Nazism and Americanism. A new paganism appeared in appeals to the laws of economics, the passions of
race and the inviolable rights of individuals. Christians were forced to suffer, to think and to unite
in new ways. Protestants reached out to one another in movements for unity; Roman Catholics struggled
to bring their church up to date. The fall of the Berlin Wall signals a decline in the power of
ideologies, while a ``new'' Christianity in the Third World and beyond emerges and grows at a rate
beyond counting. This new Christianity overshadows the inflamed Muslim voices of the time, in which it
sees both a danger and a call to mission.
\end{partoverview}
\partchapters{%
  \pch{42}{Graffiti on a Wall of Shame}{Twentieth-Century Ideologies}
  \pch{43}{Rootless Immigrants in a Sick Society}{American Evangelicals}
  \pch{44}{New Creeds for Breakfast}{The Ecumenical Movement}
  \pch{45}{The Medicine of Mercy}{Roman Catholicism: Vatican II}}
